\section*{अध्याय \ref{ch:b2:blood-circulation} --- रक्त और परिसंचरण तंत्र}

\begin{solution}{exo:b2:blood-circulation:1}
\omterm{def:b2:blood-circulation:blood}{प्लाज़्मा} \qty{55}{\%} (जल, \qty{70}{g/L} प्रोटीन, लवण,
पोषक, गैसें, हॉर्मोन); कोशिकाएँ \qty{45}{\%}। लाल कोशिकाएँ $5\times
10^{12}$/L, \qty{7}{\micro m} की चकतियाँ, ऑक्सीजन-परिवहन; श्वेत कोशिकाएँ
$7\times 10^{9}$/L, \qtyrange{10}{15}{\micro m}, प्रतिरक्षा; \omterm{def:b2:blood-circulation:blood}{बिंबाणु}
$3\times 10^{11}$/L, \qty{2}{\micro m} के टुकड़े, रक्तस्तंभन।
\end{solution}

\begin{solution}{exo:b2:blood-circulation:2}
मछली: हृदय $\to$ क्लोम $\to$ देह $\to$ हृदय, एक ही परिपथ, और देह को
\omterm{def:b2:blood-circulation:blood}{रक्त} उस नीचे दाब पर मिलता है जो क्लोमों के बाद बचा। स्तनी:
दायाँ हृदय $\to$ फेफड़े (\qty{25}{mmHg}) $\to$ बायाँ हृदय $\to$ देह
(\qty{120}{mmHg}) $\to$ दायाँ हृदय। दोहरा परिपथ देह को
फेफड़ों को उसमें डाले बिना ऊँचा दाब देता है, और दोनों
परिपथों को अलग-अलग नियंत्रित होने देता है।
\end{solution}

\begin{solution}{exo:b2:blood-circulation:3}
\omterm{def:b2:blood-circulation:vessels}{धमनी}: प्रत्यास्थ परतों तथा पेशी वाली मोटी भित्ति --- \omterm{def:b2:blood-circulation:blood}{रक्त} को
ऊँचे दाब पर ले जाती है और स्पंद चिकना करती है। \omterm{def:b2:blood-circulation:vessels}{धमनिका}: भित्ति अधिकतर
चिकनी पेशी --- प्रतिरोध तय करती और प्रवाह बाँटती है। \omterm{def:b2:blood-circulation:vessels}{केशिका}:
एक अंतःस्तरी कोशिका भर मोटी --- विनिमय। \omterm{def:b2:blood-circulation:vessels}{शिरा}: कपाटों वाली पतली,
फैलनशील भित्ति --- \omterm{def:b2:blood-circulation:blood}{रक्त} को नीचे दाब पर लौटाती है और उसका अधिकांश
संचित रखती है।
\end{solution}

\begin{solution}{exo:b2:blood-circulation:4}
$70\times 72 = \qty{5}{L/min}$, यानी प्रतिदिन \qty{7200}{L}, जबकि \omterm{def:b2:blood-circulation:blood}{रक्त}
\qty{5}{L} है: वह आयतन दिन भर में \num{1440} बार
हृदय से गुज़रता है। कोई अंग प्रतिदिन सात टन \omterm{def:b2:blood-circulation:blood}{रक्त} बना या खपा नहीं
सकता; वही \omterm{def:b2:blood-circulation:blood}{रक्त} लौटता होगा --- यानी वह परिसंचरित होता है।
\end{solution}

\begin{solution}{exo:b2:blood-circulation:5}
$\Delta P = 8\eta LQ/\pi r^{4} = 8\times 3\times 10^{-3}\times
0.4\times 8.3\times 10^{-5}/(\pi\times 2.44\times 10^{-8}) =
\qty{10}{Pa}$, यानी \qty{0.08}{mmHg}: महाधमनी कुछ नहीं लेती; वह
दाब \omterm{def:b2:blood-circulation:vessels}{धमनिकाओं} में ख़र्च होता है।
\end{solution}

\begin{solution}{exo:b2:blood-circulation:6}
$(15/12)^{4} = 2.4$: प्रतिरोध 2.4 गुना ऊपर, प्रवाह गिरकर
\qty{41}{\%}। $(15/20)^{4} = 0.32$: प्रतिरोध गिरकर एक तिहाई, प्रवाह
3.2 गुना ऊपर।
\end{solution}

\begin{solution}{exo:b2:blood-circulation:7}
महाधमनी $83/3 = \qty{28}{cm/s}$; \omterm{def:b2:blood-circulation:vessels}{केशिकाएँ} $83/3000 =
\qty{0.028}{cm/s}$; पारगमन $0.1/0.028 = \qty{3.6}{s}$। व्यायाम में
निर्गम पाँच गुना चढ़ता है और खुली केशिका-क्षेत्रफल शायद
तीन गुना, इसलिए पारगमन-काल गिरकर लगभग \qty{2}{s} रह जाता है --- फिर भी
ऑक्सीजन के निकल जाने भर को काफ़ी।
\end{solution}

\begin{solution}{exo:b2:blood-circulation:8}
धमनीय सिरा $35 - 25 = \qty{+10}{mmHg}$ (छनन); शिरीय सिरा $15
- 25 = \qty{-10}{mmHg}$ (पुनरवशोषण)। शिरीय सिरे पर $P_c = 28$ के साथ
शुद्ध $+3$ है: तरल पूरी लंबाई भर छनता है और कुछ भी
फिर नहीं सोखा जाता --- यानी शोफ, हृद्-विफलता के वे सूजे टख़ने।
\end{solution}

\begin{solution}{exo:b2:blood-circulation:9}
$40/66000 = \qty{0.61}{mol/m^3}$; $\pi = 0.61\times 8.314\times 310 =
\qty{1570}{Pa} = \qty{11.8}{mmHg}$; आधा एल्ब्यूमिन, \qty{5.9}{mmHg}।
सोडियम केशिका-भित्ति स्वतंत्र रूप से पार करता है, इसलिए उसकी सांद्रता दोनों
ओर वही है और वह उसके आरपार कोई परासरण दाब नहीं डालता।
\end{solution}

\begin{solution}{exo:b2:blood-circulation:10}
$RC = \qty{2}{s}$: $120e^{-0.4} = \qty{80}{mmHg}$। $C = 1$: $RC = 1$,
$120e^{-0.8} = \qty{54}{mmHg}$ --- और वही आघात-आयतन किसी कड़ी महाधमनी में
प्रकुंचनी दाब अधिक चढ़ा देता है। चौड़ा \omterm{thm:b2:blood-circulation:windkessel}{स्पंद दाब}
का अर्थ है कोई ऊँचा शिखर जिसके विरुद्ध हृदय को निकालना है, यानी अधिक काम
और उसी निर्गम के लिए अधिक ऑक्सीजन।
\end{solution}

\begin{solution}{exo:b2:blood-circulation:11}
विश्राम: $90/83 = \qty{1.08}{mmHg.s/mL}$; व्यायाम: $105/417 =
\qty{0.25}{mmHg.s/mL}$, यानी चार गुना पात। चूँकि $R \propto r^{-4}$,
इसलिए $r$ $4^{1/4} = 1.41$ से चढ़ा है: \omterm{def:b2:blood-circulation:vessels}{धमनिकाएँ} औसतन
\qty{41}{\%} चौड़ी हुई हैं।
\end{solution}

\begin{solution}{exo:b2:blood-circulation:12}
संतति हर स्तर पर वेग को कुल अनुप्रस्थ काट से तय कर देती
है, और वह वृक्ष ऐसे बना है कि वह काट \omterm{def:b2:blood-circulation:vessels}{केशिकाओं} में महाधमनी से
हज़ार गुना और \omterm{def:b2:blood-circulation:vessels}{शिराओं} में फिर छोटी हो;
\omterm{thm:b2:blood-circulation:poiseuille}{प्वाज़ेई का नियम} प्रतिरोध को \omterm{def:b2:blood-circulation:vessels}{धमनिकाओं} में रखता है,
जहाँ पेशी उसे बदल सकती है, और चौड़ी वाहिकाओं को हानि से लगभग
मुक्त छोड़ देता है। किसी खुले परिसंचरण के पास न तो \omterm{def:b2:blood-circulation:blood}{रक्त} को किसी नियत जगह
धीमा करने को \omterm{def:b2:blood-circulation:vessels}{केशिकाएँ} हैं, न कोई प्रतिरोध बैठाने को वाहिकाएँ, न किसी एक
अंग को प्रवाह भेजने का कोई रास्ता --- वह केवल हिला सकता है।
\end{solution}

\begin{solution}{pb:b2:blood-circulation:1}
\textbf{1.} $70\times 72 = \qty{5.0}{L/min}$; प्रतिदिन \qty{7300}{L}।
\textbf{2.} $7300/5 \approx 1450$ बार।
\textbf{3.} मिनट भर में 72 धड़कनों पर \qty{57}{g} का एक चौथाई: \qty{14}{g}
$\times 4320 = \qty{62}{kg}$ प्रति घंटा --- यानी लगभग किसी आदमी का भार।
\textbf{4.} कोई भोजन हर घंटे किसी आदमी के भार भर \omterm{def:b2:blood-circulation:blood}{रक्त} दे नहीं सकता, और
कोई ऊतक खपा नहीं सकता; एकमात्र निकास यह है कि वही \omterm{def:b2:blood-circulation:blood}{रक्त}
हृदय को लौटता है।
\textbf{5.} बाँह पर बँधी कोई डोरी उसके नीचे की \omterm{def:b2:blood-circulation:vessels}{शिराएँ} फुला देती है, ऊपर
की नहीं, इसलिए \omterm{def:b2:blood-circulation:vessels}{शिराओं} में \omterm{def:b2:blood-circulation:blood}{रक्त} हृदय की ओर चलता है;
किसी \omterm{def:b2:blood-circulation:vessels}{शिरा} में \omterm{def:b2:blood-circulation:blood}{रक्त} को हाथ की ओर दबाने पर वह कपाटों पर रुक जाता है और
पीछे धकेला नहीं जा सकता --- वे कपाट प्रवाह केवल हृदय की ओर आने देते हैं।
\textbf{6.} वे \omterm{def:b2:blood-circulation:vessels}{केशिकाएँ} जो धमनियों को \omterm{def:b2:blood-circulation:vessels}{शिराओं} से जोड़ती हैं; माल्पीगी ने
उन्हें मेंढक के फेफड़े में 1661 में देखा।
\textbf{7.} क्षेत्रफल $\pi\times 1.25^{2} = \qty{4.9}{cm^2}$; $83/4.9 =
\qty{17}{cm/s}$।
\textbf{8.} $8\times 3\times 10^{-3}\times 0.4\times 8.3\times
10^{-5}/(\pi\times 2.44\times 10^{-8}) = \qty{10}{Pa}$, \qty{0.08}{mmHg}।
\textbf{9.} $R_{1} = 8\times 3\times 10^{-3}\times 10^{-3}/(\pi\times
5.06\times 10^{-20}) = \qty{1.5e14}{Pa.s/m^3}$; समांतर में,
$1.5\times 10^{14}/3\times 10^{6} = \qty{5.0e7}{Pa.s/m^3}$।
\textbf{10.} $8.3\times 10^{-5}\times 5.0\times 10^{7} = \qty{4200}{Pa}
= \qty{31}{mmHg}$।
\textbf{11.} खुली \omterm{def:b2:blood-circulation:vessels}{केशिकाएँ} $10^{10}$, हर एक $\pi(3\times 10^{-6})^{2}
= \qty{2.8e-11}{m^2}$: \qty{0.28}{m^2} $= \qty{2800}{cm^2}$; $v =
83/2800 = \qty{0.03}{cm/s}$; पारगमन $0.1/0.03 = \qty{3.3}{s}$।
\textbf{12.} प्रतिरोध $\times(1/0.9)^{4} = 1.52$: \qty{47}{mmHg};
हृदय को धमनीय दाब \qty{16}{mmHg} चढ़ाना पड़ेगा वरना
निर्गम एक तिहाई गिर जाएगा।
\textbf{13.} $+10$, $-10$ और $0$ mmHg।
\textbf{14.} छनन $10\times 2 = \qty{20}{L}$ प्रतिदिन; पुनरवशोषण
$8.5\times 2 = \qty{17}{L}$।
\textbf{15.} प्रतिदिन \qty{3}{L} लसीका।
\textbf{16.} $\pi_c = \qty{12.5}{mmHg}$: धमनीय सिरे पर शुद्ध $+22.5$
और शिरीय सिरे पर $+2.5$ --- यानी हर जगह छनन, कोई
पुनरवशोषण नहीं: शोफ।
\textbf{17.} शुद्ध $+10$ और $+3$, औसत $+6.5$: पूरी \omterm{def:b2:blood-circulation:vessels}{केशिका} भर
छनन, यानी प्रतिदिन कोई \qty{26}{L}; लसीका वाहिकाएँ उसे ढो नहीं
सकतीं और वह तरल ऊतकों में जमा होता है, पहले पैरों में।
\textbf{18.} $4\times 10^{10}\times 2\pi\times 3\times 10^{-6}\times
10^{-3} = \qty{750}{m^2}$; $t = x^{2}/2D = (2\times 10^{-5})^{2}/
(2\times 10^{-9}) = \qty{0.2}{s}$।
\textbf{19.} $(93 - 3)/83 = \qty{1.08}{mmHg.s/mL}$।
\textbf{20.} $RC = \qty{2.2}{s}$; $120e^{-0.8/2.2} = 120\times 0.69 =
\qty{83}{mmHg}$।
\textbf{21.} $RC = \qty{1.1}{s}$; $120e^{-0.73} = \qty{58}{mmHg}$:
\omterm{thm:b2:blood-circulation:windkessel}{स्पंद दाब} 37 से चौड़ा होकर \qty{62}{mmHg} हो जाता है।
\textbf{22.} $120e^{-0.3/2.2} = \qty{105}{mmHg}$: किसी तेज़ दर पर वह
दाब धड़कनों के बीच बहुत कम गिरता है।
\textbf{23.} $C\Delta P = 2\times 20 = \qty{40}{mL}$ संचित; बाक़ी
\qty{30}{mL} प्रकुंचन के दौरान ही \omterm{def:b2:blood-circulation:vessels}{धमनिकाओं} से बह जाते हैं।
\textbf{24.} सिकुड़ता निलय प्रकुंचन में अपनी ही वाहिकाएँ भींच देता
है, इसलिए हृद्-पेशी अनुशिथिलन में सींची जाती है, उसी
दाब से जो सिकुड़ती महाधमनी बनाए रखती है; और कोई कड़ी महाधमनी, जो
अनुशिथिलनी दाब को ढहने दे, हृदय को धड़कनों के बीच भूखा रखती है।
\textbf{25.} निर्गम \qty{5}{L/min}; धमनिका-पात \qty{31}{mmHg};
लसीका को प्रतिदिन शुद्ध छनित \qty{3}{L}; अनुशिथिलनी \qty{83}{mmHg}
$C = 2$ के लिए और \qty{58}{mmHg} $C = 1$ के लिए।
\end{solution}
